% !TEX root = ../main.tex
\chapter{R and RStudio Setup}\label{app:rsetup}
\begin{paracol}{2}

Several case studies in both courses are distributed as runnable R code rather than only as PDF
write-ups: the 2013 case studies (\cref{ch:regression}) ship as plain \file{.r} scripts, and the 2024
case studies (\cref{ch:markets,ch:linalg,ch:volatility}) ship as RStudio projects with R Markdown
(\file{.Rmd}) files. This appendix collects, from the actual install scripts and instruction files
found in the course materials, what is needed to run them and which chapter each file belongs to. It
is a practical reference, not a tutorial on R itself.

\section{2013 case studies: plain R scripts}

The 18.S096 (2013) case studies used in \cref{ch:regression} come with a small family of files whose
names are hash-prefixed in \file{static\_resources/} (the hash is omitted below, as in
\cref{app:materials}): a library-install script, one or two data/plotting scripts per case study, and
a plain-text instructions file.

\begin{coursenote}[Install script pattern]
\file{fm\_casestudy\_0\_InstallOrLoadLibraries.r} (used for Case Study 1) and
\file{fm\_casestudy\_0\_InstallOrLoadLibraries Case\_Study\_2.r} (Case Study 2) both follow the same
pattern: a boolean flag \texttt{ind.install0} guards a block of \texttt{install.packages(...)} calls,
followed by \texttt{library(...)} calls that load the same packages. Set \texttt{ind.install0 <- TRUE}
the first time the script runs on a machine (or to update the packages), then switch it back to
\texttt{FALSE} on later runs so the script only loads what is already installed. The packages required
are:
\begin{center}
\texttt{quantmod}, \texttt{tseries}, \texttt{vars}, \texttt{fxregime}, \texttt{moments},
\texttt{rugarch}, \texttt{xts}
\end{center}
(\texttt{xts} appears only in the Case Study 1 variant of the install script, and only in its
\texttt{install.packages} block, not among the \texttt{library} calls; the Case Study 2 variant installs and
loads the other six.)
\end{coursenote}

The two \file{R\_Rstudio\_Instructions\_CaseStudy1/2.txt} files give the following steps (paraphrased;
both files agree except for the file lists in step 2):
\begin{enumerate}[1.]
\item Install R and RStudio (\url{http://www.r-project.org/} and
  \url{http://www.rstudio.com/ide/download/desktop}).
\item Create a project directory (e.g.\ \file{fm\_casestudy1} or \file{fm\_casestudy2}) and copy into
  it the case study's \file{.r} script files (for Case Study 2, also the \file{fred\_fxrates.txt} data
  file).
\item In RStudio: Project $\to$ Create Project $\to$ Create Project from Existing Directory, and select
  that directory.
\item Open the script files, listed in the lower-right Files pane, in the order given by the
  instructions file.
\item Source (or run with Ctrl+Enter) the install/load script first, with \texttt{ind.install0 <-
  TRUE} on the very first run; then source the data-download/plotting script; then the case study's
  \texttt{\_rcode.r} script, which reproduces the R commands shown in the case study's PDF.
\item Restart RStudio between steps if prompted (e.g.\ after installing packages), answering \texttt{n}
  to "save workspace" -- the scripts reload the data explicitly and do not rely on a saved
  \file{.RData} workspace.
\end{enumerate}

\begin{center}\footnotesize
\begin{tabular}{@{}p{0.3\linewidth}p{0.62\linewidth}@{}}
\toprule
\textbf{Case study} & \textbf{Files} \\
\midrule
Case Study 1 (regression, FX \& rates) & \file{fm\_casestudy\_0\_InstallOrLoadLibraries.r},
  \file{fm\_casestudy\_1\_0.r}, \file{fm\_casestudy\_1\_rcode.r} \\
Case Study 2 (FX regime) & \file{fm\_casestudy\_0\_InstallOrLoadLibraries Case\_Study\_2.r},
  \file{fm\_casestudy\_fx\_1.r}, \file{fm\_casestudy\_2\_rcode.r} \\
\bottomrule
\end{tabular}
\end{center}

Both case studies are used in \cref{ch:regression}; see \cref{app:materials} for the exact
correspondence with the PDF write-ups (\file{CaseStudy1.pdf}, \file{CaseStudy2.pdf}).

\section{2024 case studies: RStudio projects}

The 18.642 (2024) case studies distribute a full RStudio project per case study (a \file{.Rproj} file,
one or more \file{.Rmd} files, and any data the case study needs as an \file{.RData} or \file{.csv}
file), rather than loose scripts. Opening the \file{.Rproj} file in RStudio sets the working directory
and loads the project automatically; there is no separate install-libraries script in these archives,
so the required packages must be installed with \texttt{install.packages(...)} once, before knitting
the \file{.Rmd} file, if not already present.

\begin{center}\footnotesize
\begin{tabular}{@{}p{0.36\linewidth}p{0.13\linewidth}p{0.42\linewidth}@{}}
\toprule
\textbf{Archive / project} & \textbf{Used in} & \textbf{R packages (\texttt{library(...)})} \\
\midrule
\texttt{\small mit18\_642\_f24\_\allowbreak rstudio.zip} (\file{FM\_Intro1.Rmd}) & \cref{ch:markets} &
  \texttt{tidyquant}, \texttt{tidyverse}, \texttt{ggplot2} \\
\texttt{\small mit18\_642\_f24\_\allowbreak equal\_weighted\_\allowbreak portfolios.zip} & \cref{ch:linalg} &
  \texttt{tidyquant}, \texttt{tidyverse}, \texttt{ggplot2} \\
\texttt{\small mit18\_642\_f24\_\allowbreak etf\_casestudy.zip} & \cref{ch:regression} &
  \texttt{tidyquant}, \texttt{tidyverse}, \texttt{car}, \texttt{broom} \\
\texttt{\small mit18\_642\_f24\_\allowbreak rstudio\_pset5.zip} (ARKK/S\&P 500 volatility) & \cref{ch:volatility} &
  \texttt{tidyquant}, \texttt{tidyverse}, \texttt{fpp2}, \texttt{ggfortify}, \texttt{glmnet} \\
\bottomrule
\end{tabular}
\end{center}

(This is the union of \texttt{library(...)} calls actually found across the \file{.Rmd} files in these
four archives; a given \file{.Rmd} may only use a subset of the packages listed for its row. No
combined install script for the 2024 material was found in the static resources -- unlike the 2013
scripts, package installation here is left implicit and must be done once per package with
\texttt{install.packages("<name>")} before the first knit.)

\begin{coursenote}
Two of these are financial-data packages worth calling out: \texttt{tidyquant} and
\texttt{quantmod} both pull historical price/rate series from the internet (typically Yahoo Finance or
FRED), so an internet connection is needed the first time a script or \file{.Rmd} runs, even though the
2013 scripts also cache the result in an \file{.RData} workspace for later offline reruns.
\end{coursenote}

\section{Try it yourself}

\begin{itemize}
\item Run \file{FM\_Intro1.Rmd} (\cref{ch:markets}) end to end in RStudio with "Knit" once the
  packages above are installed, and compare the output to \file{FM\_Intro1.pdf} in the same archive.
\item Change \texttt{ind.install0} back and forth in
  \texttt{\small fm\_casestudy\_0\_\allowbreak InstallOrLoadLibraries.r} and confirm that with it set
  to \texttt{FALSE} the script still runs (using the already-installed packages) but skips the
  \texttt{install.packages} calls.
\end{itemize}
